% !PRINTED 101
\weberpage{101}
Aus der Zusammensetzung
\weq{7}{\begin{pmatrix}1, & 0\\0, & n\end{pmatrix}\begin{pmatrix}1, & 0\\0, & m\end{pmatrix} = \begin{pmatrix}1, & 0\\0, & mn\end{pmatrix}, \quad \begin{pmatrix}1, & 0\\0, & n\end{pmatrix}\begin{pmatrix}n, & 0\\0, & 1\end{pmatrix} = \begin{pmatrix}n, & 0\\0, & n\end{pmatrix}}
können wir noch weiter schließen, daß sich \gesperrt{jede Transformation
vom Grade $n$ aus solchen zusammensetzen läßt, deren
Grad eine Primzahl ist.} Zerlegt man $n = pq$ in zwei Faktoren
 $p$ und $q$, die zueinander relativ prim sind, so ergibt sich,
indem man die Zahlen $\beta, \delta$ aus
\[
p\delta - q\beta = 1
\]
bestimmt, die Zusammensetzung
\[
\begin{pmatrix}p, & \beta\\q, & \delta\end{pmatrix}\begin{pmatrix}1, & 0\\0, & n\end{pmatrix}\begin{pmatrix}p\delta, & -q\beta\\-1, & 1\end{pmatrix} = \begin{pmatrix}p, & 0\\0, & q\end{pmatrix},
\]
woraus zu ersehen ist, daß man statt der Haupttransformation
auch jede dieser Transformationen $\begin{pmatrix}p, & 0\\0, & q\end{pmatrix}$ zur Ableitung aller
anderen benutzen kann.

\begin{center}
\textbf{§ 30. \gesperrt{Die linearen Fundamentaltransformationen.}}
\end{center}

Die ganze Gruppe $\mathfrak{L}$ der linearen Transformationen läßt sich
durch Wiederholung von zweien unter ihnen, die wir die \gesperrt{linearen
Fundamentaltransformationen} nennen, ableiten.
Ist
\[
L = \begin{pmatrix}\alpha, & \beta\\\gamma, & \delta\end{pmatrix}, \quad \alpha\delta - \beta\gamma = 1
\]
eine beliebige lineare Transformation, so ist
\weq{1}{\begin{pmatrix}\alpha, & \beta\\\gamma, & \delta\end{pmatrix}\begin{pmatrix}\delta, & -\beta\\-\gamma, & \alpha\end{pmatrix} = \begin{pmatrix}1, & 0\\0, & 1\end{pmatrix},}
und da die identische Transformation die Einheit in der Gruppe $L$
ist, % [sic] Printed L; the group is otherwise denoted $\mathfrak{L}$.
so ist
\[
L^{-1} = \begin{pmatrix}\delta, & -\beta\\-\gamma, & \alpha\end{pmatrix}
\]
die zu $L$ \gesperrt{reziproke Transformation}.

Wir bezeichnen durch die Potenz $L^{\pm m}$ das, was durch $m$ malige
Wiederholung von $L$ oder $L^{-1}$ entsteht, und wollen nun nachweisen,
daß sich durch die Potenzen der \gesperrt{Fundamentaltransformationen}

\weq{2}{A = \begin{pmatrix}1, & 0\\1, & 1\end{pmatrix}, \quad B = \begin{pmatrix}0, & 1\\-1, & 0\end{pmatrix}}
%CONT>>
